% !TEX program = lualatex
% Документація до класу dstu4163. Компіляція: latexmk (LuaLaTeX) або xelatex.
% Перед компіляцією зберіть examples/ і templates/ (make) — їхні PDF
% вставляються в розділ «Галерея прикладів».
\documentclass[11pt,a4paper]{article}

\usepackage[a4paper,left=25mm,right=25mm,top=25mm,bottom=25mm]{geometry}
\usepackage{fontspec}
\IfFontExistsTF{Times New Roman}
  {\setmainfont{Times New Roman}}
  {\IfFontExistsTF{Liberation Serif}{\setmainfont{Liberation Serif}}{\setmainfont{TeX Gyre Termes}}}
\IfFontExistsTF{Menlo}
  {\setmonofont[Scale=0.85]{Menlo}}
  {\IfFontExistsTF{DejaVu Sans Mono}
     {\setmonofont[Scale=0.85]{DejaVu Sans Mono}}
     {\setmonofont{Liberation Mono}}}
\usepackage{babel}
\babelprovide[main,import]{ukrainian}

\usepackage{graphicx}
\usepackage{array}
\usepackage{booktabs}
\usepackage{longtable}
\usepackage{fancyvrb}
\usepackage{xcolor}
\usepackage{tikz}
\usetikzlibrary{arrows.meta}
\usepackage{enumitem}
\usepackage[font=small,labelfont=bf]{caption}
\usepackage[hidelinks]{hyperref}
\hypersetup{
  pdftitle={dstu4163 — українські листи за ДСТУ 4163:2020},
  pdfsubject={Документація до класу LaTeX},
  pdfkeywords={ДСТУ 4163:2020, LaTeX, службовий лист, діловодство}
}

\definecolor{accent}{HTML}{1F4E79}
\definecolor{codebg}{HTML}{F4F6F8}

\setlength\parindent{0pt}
\setlength\parskip{0.5\baselineskip}
\setlist{nosep,leftmargin=*}
\setcounter{tocdepth}{2}
\setlength\emergencystretch{2em}

\newcommand\pkg[1]{\textsf{#1}}
\newcommand\cls{\pkg{dstu4163}}
\newcommand\cs[1]{\texttt{\char`\\#1}}
\newcommand\key[1]{\texttt{#1}}
\newcommand\val[1]{\texttt{#1}}
\newcommand\meta[1]{$\langle$\textit{#1}$\rangle$}
\newcommand\req[1]{\textsc{р}\,#1}
\newcommand\cl[1]{п.\,#1}

\DefineVerbatimEnvironment{code}{Verbatim}
  {fontsize=\small,frame=leftline,framerule=1.2pt,rulecolor=\color{accent},
   framesep=3mm,xleftmargin=2mm}

\newcolumntype{L}[1]{>{\raggedright\arraybackslash}p{#1}}
\setlength\tabcolsep{4pt}

\graphicspath{{../examples/}{../templates/}}
\newcommand\thumb[2][0.3]{%
  {\setlength\fboxsep{0pt}\fbox{\includegraphics[width=\dimexpr#1\linewidth-0.8pt\relax,page=1]{#2}}}}

\begin{document}

% =============================================================================
\begin{titlepage}
  \centering
  \vspace*{30mm}
  {\Huge\bfseries\color{accent} dstu4163\par}
  \vspace{6mm}
  {\Large Українські службові та ділові листи\\ за ДСТУ 4163:2020 у \LaTeX\par}
  \vspace{10mm}
  {\large Версія 1.1 \quad·\quad 18.11.2026\par}
  \vspace{4mm}
  {\large Dmytro\par}
  \vfill
  \begin{minipage}{0.8\linewidth}
    \small
    Клас \cls{} верстає листи відповідно до ДСТУ 4163:2020 «Уніфікована система
    організаційно-розпорядчої документації. Вимоги до оформлення документів»:
    кутовий і поздовжній бланки, адресат, заголовок, додатки, підпис,
    відомості про виконавця, нумерація сторінок. Розміри полів і відступів
    задано стандартом; порожні реквізити не друкуються, тож один шаблон
    придатний для службового листа, супровідного й гарантійного листів,
    особистого ділового листа, звернення громадянина та запиту на публічну
    інформацію.
  \end{minipage}
  \vspace{20mm}
\end{titlepage}

\tableofcontents
\newpage

% =============================================================================
\section{Вступ}

\subsection{Призначення}

ДСТУ 4163:2020 затверджено наказом ДП «УкрНДНЦ» від 01.07.2020 № 144; він
чинний з 01.09.2021 і замінив ДСТУ 4163-2003. Стандарт поширюється на
організаційно-розпорядчі документи незалежно від носія — паперові та
електронні. Для юридичних осіб його застосування обов'язкове відповідно до
Правил організації діловодства та архівного зберігання документів (наказ
Мін'юсту від 18.06.2015 № 1000/5); фізичні особи-підприємці та громадяни
можуть застосовувати його добровільно.

Клас \cls{} реалізує ті вимоги стандарту, що стосуються листів:

\begin{itemize}
  \item поля сторінки, гарнітуру, кеглі й інтервали (\cl{6.2}, \cl{7.2}, \cl{7.3});
  \item кутовий і поздовжній бланки листа зі зоною реквізитів 73\,мм
        (\cl{6.6}, \cl{6.7}, додаток~А);
  \item відступи від лівого поля: абзац 10\,мм, адресат 90\,мм, гриф 100\,мм,
        розшифрування підпису 125\,мм (\cl{7.7});
  \item реквізити 01--15, 19--23, 27 з правилами їхнього розміщення;
  \item нумерацію сторінок з другої (\cl{7.10});
  \item автоматичне вилучення порожніх реквізитів (розд.~\ref{sec:dropout}).
\end{itemize}

\subsection{Склад пакета}

\begin{longtable}{@{}L{58mm}L{96mm}@{}}
\toprule
Файл & Призначення \\
\midrule
\endhead
\texttt{dstu4163.cls} & клас \\
\texttt{README.md} & короткий опис (англ.) \\
\texttt{Makefile} & збірка, встановлення, очищення \\
\texttt{doc/dstu4163-doc.pdf} & цей посібник \\
\texttt{templates/lyst-porozhnii.tex} & порожній шаблон з примітками до кожного поля \\
\texttt{templates/lyst-povnyi.tex} & повний шаблон: задіяно всі можливості \\
\texttt{examples/00-…07-…} & вісім типових листів (розд.~\ref{sec:gallery}) \\
\texttt{*/latexmkrc} & налаштування \texttt{latexmk} для кожної теки \\
\bottomrule
\end{longtable}

% =============================================================================
\section{Встановлення та компіляція}

\subsection{Вимоги}

\begin{itemize}
  \item TeX Live 2021 або новіший (MacTeX, MiKTeX); рушій \textbf{LuaLaTeX}
        (рекомендовано) або \textbf{XeLaTeX}. pdf\LaTeX{} не підтримується:
        клас використовує \pkg{fontspec} і системні шрифти.
  \item Пакети: \pkg{fontspec}, \pkg{babel} (українська через
        \cs{babelprovide}), \pkg{geometry}, \pkg{graphicx}, \pkg{iftex}.
  \item Шрифт Times New Roman (\cl{7.2}). Якщо його немає, клас бере
        Liberation Serif, а далі TeX Gyre Termes — метрично сумісні гарнітури.
  \item Українські шаблони переносів \texttt{hyph-uk} (колекція
        \texttt{collection-langcyrillic}; у повних MacTeX і TeX Live вже є).
\end{itemize}

\subsection{Варіанти встановлення}

\textbf{Поруч із документом.} Покладіть \texttt{dstu4163.cls} у теку з
\texttt{.tex}-файлом.

\textbf{Для всіх документів користувача.} Скопіюйте клас у \texttt{TEXMFHOME}:
\begin{code}
make install        # копіює в $(kpsewhich -var-value TEXMFHOME)/tex/latex/dstu4163
make uninstall      # вилучає
\end{code}
На macOS це \texttt{\textasciitilde/Library/texmf/tex/latex/dstu4163/}. Після
встановлення клас знаходять усі редактори (TeXShop, VS~Code, Overleaf~— ні:
туди файл треба завантажити разом із документом).

\textbf{У межах пакета.} Файли \texttt{latexmkrc} у теках \texttt{examples},
\texttt{templates} і \texttt{doc} додають батьківську теку до шляху пошуку,
тому \texttt{latexmk} там працює без встановлення.

\subsection{Збирання}

\begin{code}
make                # приклади, шаблони, документація (LuaLaTeX)
make ENGINE=xelatex # те саме через XeLaTeX
make clean          # вилучити допоміжні файли
make dist           # архів dstu4163.zip

cd templates && latexmk lyst-povnyi.tex        # один документ
\end{code}

\subsection{LuaLaTeX чи XeLaTeX}

Для класу рушії рівноцінні: розміри, відступи й кеглі однакові. Розриви
рядків можуть дещо відрізнятися через різне формування тексту. LuaLaTeX
активно розвивається, підтримує розтягування шрифту в \pkg{microtype} і
завантажує шаблони переносів на льоту; XeLaTeX швидший. Типово
використовується LuaLaTeX; щоб зібрати через XeLaTeX, виконайте
\texttt{latexmk -xelatex}.

% =============================================================================
\section{Швидкий старт}

\begin{code}
\documentclass[14pt]{dstu4163}
\dstusetup{
  name    = {Товариство з обмеженою відповідальністю «Альфа»},
  address = {вул. Хрещатик, 1,\\ м. Київ, 01001},
  phone   = {(044) 000-00-00},
  email   = {office@alfa.example},
  code    = {12345678},
  date    = {18.11.2026},
  index   = {01-12/345},
}
\addressee{Директору ТОВ «Бета»\\ Роману \surname{Зуєву}}

\begin{document}
\heading{Про співпрацю}
\salutation{Шановний пане Романе!}

Пропонуємо обговорити умови спільного дослідницького проєкту.

\signature{Директор}{Ігор \surname{Шевчук}}
\executor{Мельник Ольга 044 000 00 01}
\end{document}
\end{code}

Бланк і адресат друкуються автоматично перед першою з команд \cs{heading},
\cs{salutation}; якщо їх немає (наприклад, у заяві), викличте
\cs{makeletterhead}. Відомості про виконавця завжди потрапляють у лівий
нижній кут останньої сторінки.

% =============================================================================
\section{Опції класу}

\begin{longtable}{@{}L{32mm}L{120mm}@{}}
\toprule
Опція & Дія \\
\midrule
\endhead
\val{12pt}, \val{13pt}, \val{14pt} & Кегль основного тексту (\cl{7.2}: 12--14\,пт).
  Типово \val{14pt}. Довідкові дані друкуються на 2\,пт дрібніше, виконавець~— 10\,пт,
  назва виду документа~— на 2\,пт більше. \\
\val{single}, \val{onehalf} & Міжрядковий інтервал \emph{тексту} листа (\cl{7.3}:
  1--1,5). Типово \val{single}. Багаторядкові реквізити завжди мають одинарний
  інтервал. «Одинарний»~= 1,15 кегля, як у Word для Times New Roman. \\
\bottomrule
\end{longtable}

% =============================================================================
\section{Налаштування: \cs{dstusetup}}

\cs{dstusetup}\verb|{|\meta{ключ}\,=\,\meta{значення}, \dots\verb|}| задає
реквізити й параметри. Команду можна викликати кілька разів; порожні рядки
між ключами дозволені. Значення з комами чи знаками \verb|=| беріть у фігурні
дужки. Колонка «Р» містить номер реквізиту за ДСТУ~4163:2020.

\subsection{Бланк і автор}

\begin{longtable}{@{}L{33mm}L{28mm}L{10.5mm}L{75.5mm}@{}}
\toprule
Ключ & Значення (\textit{типово}) & Р & Опис \\
\midrule
\endhead
\key{blank} & \val{corner} | \val{long} | \val{none} \newline \textit{corner} & --- &
  Кутовий бланк, поздовжній бланк або аркуш без бланка (\cl{6.6}). У режимі
  \val{none} друкуються лише гриф, адресати, заявник і, за наявності, дата з
  індексом. \\
\key{sender} & \val{organisation} | \val{person} \newline \textit{organisation} & --- &
  Юрособа: найменування великими літерами (\cl{7.4}), «Код ЄДРПОУ» перед кодом,
  у режимі \key{form} — поле «№». Фізособа: ім'я друкується як задано, без «№». \\
\key{align} & \val{center} | \val{flush} \newline \textit{center} & --- &
  Центроване чи прапорове розташування реквізитів кутового бланка (\cl{6.7}). \\
\key{form} & \val{true} | \val{false} \newline \textit{false} & 10--12 &
  Бланк для друку: порожні дата, «№» і «На № … від …» стають лініями для
  заповнення від руки. Без \key{form} порожні реквізити зникають. \\
\key{arms}, \key{emblem} & код \TeX & 01, 02 &
  Герб (лише для органів влади) або емблема / знак для товарів і послуг.
  Рекомендована висота — до 17\,мм; емблема не замінює найменування (\cl{5.2}).
  Значення беріть у дужки: \newline\verb|{\includegraphics[height=17mm]{logo}}|. \\
\key{parent} & текст & 03 & Найменування юрособи вищого рівня (дрібнішим кеглем, \cl{7.2}). \\
\key{name} & текст & 04 & Повне найменування за статутом або ім'я фізособи. \\
\key{short-name} & текст & 04 & Скорочене найменування, якщо його зафіксовано у статуті
  (з дужками чи без, \cl{5.4}). \\
\key{unit} & текст & 05 & Структурний підрозділ, якщо його керівник має право підпису. \\
\key{address} & текст & 06 & Поштова адреса; рядки — через \verb|\\|. \\
\key{phone}, \key{email}, \key{web} & текст & 06 & Друкуються як «тел. …»,
  «e-mail: …», адреса сайту. Код міста — у дужках. \\
\key{extra} & текст & 06 & Довільний додатковий рядок (рахунок, факс). \\
\key{code} & цифри & 08 & Код ЄДРПОУ. \\
\bottomrule
\end{longtable}

\subsection{Реквізити документа}

\begin{longtable}{@{}L{33mm}L{28mm}L{10.5mm}L{75.5mm}@{}}
\toprule
Ключ & Значення & Р & Опис \\
\midrule
\endhead
\key{date} & текст & 10 & Дата: \val{18.11.2026}, \cs{dstutoday},
  \verb|\dstudateverbal{18}{11}{2026}| (\cl{5.10}). \\
\key{index} & текст & 11 & Реєстраційний індекс; складники — через правобічну
  похилу риску: \val{01-12/345} (\cl{5.11}). \\
\key{ref-index}, \key{ref-date} & текст & 12 & Індекс і дата листа, на який дають
  відповідь (\cl{5.12}). \\
\key{access} & текст & 14 & Гриф обмеження доступу; праворуч, над адресатом,
  з відступом 100\,мм (\cl{5.14}). \\
\key{doctype} & текст & 09 & Назва виду документа: великі літери, розріджено,
  на 2\,пт більше. У звичайних листах не зазначається (\cl{4.4}); потрібна для
  гарантійного листа, заяви, запиту. \\
\key{doctype-subtitle} & текст & 09 & Рядок під назвою (напр., «на отримання
  публічної інформації»). \\
\bottomrule
\end{longtable}

\subsection{Оформлення}

Типові значення відповідають ДСТУ~4163:2020; змінюйте їх, лише якщо цього
вимагає інструкція з діловодства вашої установи.

\begin{longtable}{@{}L{32mm}L{22mm}L{98mm}@{}}
\toprule
Ключ & Типово & Опис \\
\midrule
\endhead
\key{heading-bold} & \val{false} & Напівжирний заголовок до тексту (\cl{7.2} дозволяє). \\
\key{page-numbers} & \val{true} & Номери сторінок з другої, посередині верхнього поля (\cl{7.10}). \\
\key{text-spacing} & \val{1} & Інтервал тексту; \val{1.5} — те саме, що \val{onehalf}. \\
\key{gap} & \val{1} & Порожніх рядків між реквізитами (\cl{7.3}: 1,5--3 інтервали
  між базовими лініями). \\
\key{signature-gap} & \val{2} & Порожніх рядків перед першим підписом. \\
\key{zone-width} & \val{73mm} & Ширина зони кутового бланка, заголовка, посади
  в підписі (\cl{7.6}: 73\,мм, 28 знаків). \\
\key{addressee-offset} & \val{90mm} & Відступ адресата від лівого поля (\cl{7.7}). \\
\key{addressee-width} & \val{73mm} & Ширина адресата (\cl{7.6}). \\
\key{access-offset} & \val{100mm} & Відступ грифа обмеження доступу (\cl{7.7}). \\
\key{signature-offset} & \val{125mm} & Відступ розшифрування підпису (\cl{7.7}). \\
\key{paragraph-indent} & \val{10mm} & Абзацний відступ (\cl{7.7}). \\
\bottomrule
\end{longtable}

\subsection{Мітки}

Для листів іноземною мовою (\cl{5.20}) або нестандартних формулювань:

\begin{longtable}{@{}L{36mm}L{34mm}L{82mm}@{}}
\toprule
Ключ & Типово & Де друкується \\
\midrule
\endhead
\key{label-no} & № & рядок дати \\
\key{label-ref}, \key{label-from} & На №, від & посилання на вхідний лист \\
\key{label-phone}, \key{label-email} & тел., e-mail: & довідкові дані \\
\key{label-code} & Код ЄДРПОУ & код юрособи \\
\key{label-attachment} & Додаток: & \cs{attachment} \\
\key{label-attachments} & Додатки: & середовище \texttt{attachments} \\
\key{label-seal} & М. П. & \cs{signature}\verb|[seal]| \\
\bottomrule
\end{longtable}

% =============================================================================
\section{Команди}

Будь-яка команда з порожнім аргументом нічого не друкує й не додає відступів.

\begin{description}[style=nextline,leftmargin=8mm,font=\normalfont]
\item[\cs{addressee}\texttt{\{}\meta{адресат}\texttt{\}}]
  Реквізит 15. До чотирьох викликів; за більшої кількості клас попереджає:
  потрібен список розсилки (\cl{5.15}). Рядки розділяйте \verb|\\|, складові
  частини (організація, посада й ім'я, адреса) — порожнім рядком; між ними
  буде інтервал 1,5 (\cl{7.3}).
\item[\cs{applicant}\texttt{\{}\meta{заявник}\texttt{\}}]
  Блок заявника чи запитувача під адресатом — для заяв, скарг, запитів.
\item[\cs{makeletterhead}]
  Друкує бланк, гриф, адресатів і заявника. Викликається автоматично
  командами \cs{heading} і \cs{salutation}; повторний виклик ігнорується.
  Якщо всі реквізити бланка порожні, нічого не друкує.
\item[\cs{heading}\texttt{\{}\meta{заголовок}\texttt{\}}]
  Реквізит 19: без відступу, шириною до 73\,мм, без крапки в кінці (\cl{5.19}, \cl{7.6}).
\item[\cs{salutation}\texttt{\{}\meta{звертання}\texttt{\}}]
  Звертання по центру, перед текстом.
\item[\cs{attachment}\texttt{[}\meta{мітка}\texttt{]\{}\meta{текст}\texttt{\}}]
  Реквізит 21 з одним додатком: «Додаток: на 3 арк. в 1 прим.». Наступні рядки
  вирівнюються під текстом.
\item[\texttt{\char`\\begin\{attachments\}[}\meta{мітка}\texttt{]} \dots\ \cs{item} \dots\ \texttt{\char`\\end\{attachments\}}]
  Нумерований перелік додатків: «Додатки: 1. … 2. …».
\item[\cs{closing}\texttt{\{}\meta{формула}\texttt{\}}]
  Заключна формула ввічливості («З повагою») перед підписом. Не є реквізитом
  стандарту.
\item[\cs{signature}\texttt{[}\meta{ключі}\texttt{]\{}\meta{посада}\texttt{\}\{}\meta{ім'я}\texttt{\}}]
  Реквізит 22. Посада — без відступу, до 73\,мм; власне ім'я та ПРІЗВИЩЕ —
  з відступом 125\,мм на рівні останнього рядка посади (\cl{7.5}). Якщо ім'я
  довше за решту рядка, воно вирівнюється по правому полю. Ключі:
  \key{seal} — «М. П.» під підписом; \key{date}\,=\,\meta{дата} — дата
  ліворуч, коли посаду не вказано (\cl{5.10}: у документах не на бланку дату
  ставлять ліворуч від підпису); \key{gap}\,=\,\meta{рядки} — інтервал перед
  підписом. Кілька підписів — кілька викликів поспіль, за підпорядкованістю.
\item[\cs{jointsignatures}\texttt{\{}\meta{посада}\texttt{\}\{}\meta{ім'я}\texttt{\}\{}\meta{посада}\texttt{\}\{}\meta{ім'я}\texttt{\}}]
  Підписи осіб з однаковими посадами на одному рівні (спільні документи).
\item[\cs{executor}\texttt{\{}\meta{виконавець}\texttt{\}}]
  Реквізит 27: «Прізвище Власне ім'я телефон» (\cl{5.27}), кегль 10\,пт,
  лівий нижній кут останньої сторінки. Можна викликати 2--3 рази.
\item[\cs{surname}\texttt{\{}\meta{прізвище}\texttt{\}}, \cs{fullname}\texttt{\{}\meta{ім'я}\texttt{\}\{}\meta{прізвище}\texttt{\}}]
  Прізвище великими літерами: \verb|Роману \surname{Зуєву}| → Роману ЗУЄВУ.
\item[\cs{dstudate}\texttt{\{}\meta{д}\texttt{\}\{}\meta{м}\texttt{\}\{}\meta{р}\texttt{\}}, \cs{dstudateverbal}, \cs{dstudateiso}]
  18.11.2026; 18 листопада 2026 року (із зірочкою — «р.»); 2026.11.18
  (дозволено в листуванні з іноземними партнерами).
\item[\cs{dstutoday}, \cs{dstutodayverbal}]
  Поточна дата в цифровому і словесно-цифровому форматах.
\end{description}

% =============================================================================
\section{Порожні реквізити та перемикачі}\label{sec:dropout}

Клас не вимагає видаляти невикористані рядки шаблону: порожнє значення ключа
чи порожній аргумент команди просто вимикають реквізит.

\begin{longtable}{@{}L{46mm}L{54mm}L{52mm}@{}}
\toprule
Реквізит & Якщо порожній & Із \key{form = true} \\
\midrule
\endhead
01--08 реквізити бланка & рядок не друкується & так само \\
Увесь бланк (усе порожнє) & жодного відступу зверху & так само \\
10 дата & не друкується & лінія 30\,мм \\
11 індекс («№ …») & не друкується & лінія (лише юрособа) \\
12 «На № … від …» & не друкується & лінії (юрособа, є бланк) \\
14 гриф, 09 назва виду & не друкуються & так само \\
15 адресат, заявник & не друкуються & так само \\
19 заголовок, звертання & не друкуються & так само \\
21 додатки, формула ввічливості & не друкуються & так само \\
22 підпис без посади & ім'я на 125\,мм, ліворуч порожньо або дата & так само \\
27 виконавець & не друкується & так само \\
Номери сторінок & \key{page-numbers = false} вимикає & так само \\
\bottomrule
\end{longtable}

% =============================================================================
\section{Геометрія сторінки}

Рис.~\ref{fig:layout} показує розміщення реквізитів листа на кутовому
бланку. Допуск розташування реквізитів — $\pm 2$\,мм (\cl{6.5}).

\begin{figure}[htbp]
\centering
\begin{tikzpicture}[x=0.6mm,y=-0.6mm,font=\scriptsize,
    dim/.style={{Stealth[length=1.6mm]}-{Stealth[length=1.6mm]},thin,accent},
    dl/.style={midway,fill=white,inner sep=1pt,text=accent},
    zone/.style={draw=black!45,dashed,fill=accent!6}]
  % аркуш і межі текстового поля
  \draw[thick] (0,0) rectangle (210,297);
  \draw[black!35] (30,20) rectangle (200,277);
  \node[text=black!60,align=center] at (168,10) {«2» — номер сторінки\\ (з другої, по центру)};
  % розміри у верхньому полі
  \draw[dim] (30,14) -- node[dl] {90} (120,14);
  \draw[dim] (30,6) -- node[dl] {100} (130,6);
  % зони
  \draw[zone] (30,20) rectangle (103,88);
  \node[align=center] at (66.5,50) {01–08 реквізити бланка\\[3pt] 10 дата, 11 індекс,\\ 12 «На № … від …»};
  \draw[zone] (120,20) rectangle (193,62);
  \node[align=center] at (156.5,41) {14 гриф (з 100\,мм)\\[3pt] 15 адресат};
  \draw[dim] (30,92) -- node[dl] {73} (103,92);
  \draw[dim] (120,66) -- node[dl] {73} (193,66);
  \draw[zone] (30,98) rectangle (103,110);
  \node at (66.5,104) {19 заголовок};
  \draw[zone] (30,118) rectangle (200,205);
  \foreach \y in {136,142,148,154,160} {\draw[black!30] (30,\y) -- (200,\y);}
  \draw[black!30] (40,130) -- (200,130);
  \draw[dim] (30,124) -- node[dl] {10} (40,124);
  \node[align=center,fill=accent!6] at (115,182) {20 текст};
  \draw[zone] (30,212) rectangle (200,220);
  \node at (115,216) {21 додатки};
  \draw[zone] (30,228) rectangle (103,236);
  \node at (66.5,232) {22 посада};
  \draw[zone] (155,228) rectangle (200,236);
  \node at (177.5,232) {Ім'я ПРІЗВИЩЕ};
  \draw[dim] (30,244) -- node[dl] {125} (155,244);
  \draw[zone] (30,266) rectangle (103,275);
  \node at (66.5,270.5) {27 виконавець};
  % поля
  \draw[dim] (0,286) -- node[dl] {30} (30,286);
  \draw[dim] (200,286) -- node[dl] {10} (210,286);
  \draw[dim] (218,0) -- node[right,text=accent] {20} (218,20);
  \draw[dim] (218,277) -- node[right,text=accent] {20} (218,297);
\end{tikzpicture}
\caption{Схема листа на кутовому бланку, розміри в міліметрах (масштаб довільний).}
\label{fig:layout}
\end{figure}

\begin{longtable}{@{}L{66mm}L{30mm}L{56mm}@{}}
\toprule
Параметр & Значення & Підстава \\
\midrule
\endhead
Формат & А4, 210$\times$297\,мм & \cl{6.1} \\
Поля: ліве / праве / верхнє / нижнє & 30 / 10 / 20 / 20\,мм & \cl{6.2} \\
Гарнітура, кегль тексту & Times New Roman, 12--14\,пт & \cl{7.2} \\
Довідкові дані, виконавець & 8--12\,пт & \cl{7.2} \\
Назва виду документа & 14--16\,пт, розріджено & \cl{7.2}, \cl{7.4} \\
Інтервал тексту & 1--1,5 & \cl{7.3} \\
Інтервал багаторядкових реквізитів & 1 & \cl{7.3} \\
Між частинами адресата & 1,5 & \cl{7.3} \\
Між реквізитами & 1,5--3 & \cl{7.3} \\
Макс. довжина рядка реквізиту & 73\,мм (28 знаків) & \cl{7.6} \\
Абзац / адресат / гриф / підпис & 10 / 90 / 100 / 125\,мм & \cl{7.7} \\
Нумерація сторінок & з другої, по центру вгорі & \cl{7.10} \\
\bottomrule
\end{longtable}

% =============================================================================
\section{Правила оформлення: довідка}

Клас відповідає за розміщення; зміст реквізитів формує автор. Нижче —
вимоги ДСТУ~4163:2020 до змісту реквізитів листа.

\subsection{Дата (реквізит 10, \cl{5.10})}
Цифровий спосіб: день, місяць, рік арабськими цифрами через крапку, рік —
чотирма цифрами, без крапки в кінці: \val{18.11.2026}. Словесно-цифровий:
«18 листопада 2026 року» (дозволено «р.»). У листуванні з іноземними
партнерами допустимо \val{2026.11.18}. У документах не на бланку дату
ставлять ліворуч від підпису.

\subsection{Реєстраційний індекс (реквізит 11, \cl{5.11})}
Порядковий номер, доповнений індексами (справи за номенклатурою, підрозділу
тощо); складники — через правобічну похилу риску, напр.
\val{01-12/345} або \val{123/01-10}. Дату та індекс проставляють у день підписання листа.

\subsection{Адресат (реквізит 15, \cl{5.15})}
Адресат розміщують праворуч угорі першої сторінки; адресатів може бути до
чотирьох, більше — складають список розсилки.
\begin{itemize}
  \item Юрособа без посадової особи — найменування у називному відмінку:
        «Науково-дослідний центр «Зразок»».
  \item Керівник юрособи — найменування юрособи входить до посади, усе в
        давальному відмінку: «Директору ТОВ «Бета» Роману ЗУЄВУ».
  \item Керівник підрозділу — юрособа в називному, посада й ім'я в давальному:
        «Центр «Зразок» / Завідувачу відділу / Олені КОВАЛЬ».
  \item Багатьом однорідним адресатам — узагальнено в давальному відмінку.
  \item Фізична особа — власне ім'я та ПРІЗВИЩЕ в давальному відмінку, потім
        поштова адреса.
\end{itemize}
Поштову адресу не зазначають на документах, що надсилаються органам
законодавчої й виконавчої влади та постійним кореспондентам. Формат адреси —
відповідно до Правил надання послуг поштового зв'язку.

\subsection{Заголовок і текст (реквізити 19, 20, \cl{5.19}, \cl{5.20})}
Заголовок стисло передає зміст і відповідає на питання «про що?»: «Про
надання статистичних даних». Довжина рядка — до 73\,мм; заголовок понад
150 знаків можна продовжити до правого поля; крапку в кінці не ставлять.

Текст має бути стислим, точним і без повторів. Документи складають
державною мовою; зарубіжним адресатам — за домовленістю мовою адресата чи
мовою міжнародного спілкування; до державних органів — лише українською.
Від імені юрособи текст зазвичай викладають від першої особи множини
(«повідомляємо», «просимо»), від фізособи — від першої особи однини.

\subsection{Додатки (реквізит 21, \cl{5.21})}
Відмітку оформлюють після тексту перед підписом; у супровідних листах вона
обов'язкова.
\begin{code}
\attachment{на 7 арк. у 2 прим.}
\begin{attachments}
  \item Довідка про виконання плану на 5 арк. в 1 прим.
  \item Графік ремонтних робіт на 3 арк. в 1 прим.
\end{attachments}
\attachment{лист ТОВ «Зразок» від 20.09.2023 № 595/04-12 і додаток до
  нього, всього на 20 арк. в 1 прим.}
\attachment{форми галузевої звітності в 1 прим.}        % зброшуровані
\attachment[Додатки:]{згідно з описом на 2 арк.}        % понад 10 додатків
\attachment{на 12 арк. в 1 прим. на першу адресу.}     % не всім адресатам
\end{code}

\subsection{Підпис і печатка (реквізити 22, 23, \cl{5.22}, \cl{5.23})}
Підпис містить найменування посади (на бланку — у скороченій формі),
особистий підпис і власне ім'я та ПРІЗВИЩЕ: «Директор \quad Ігор ШЕВЧУК».
Кількох посадових осіб підписують одна під одною за підпорядкованістю; осіб
з однаковими посадами — на одному рівні. «В.~о.» додають лише в разі
заміщення за наказом; прийменник «За» та похилу риску перед посадою не
ставлять. Відбиток печатки має охоплювати кілька останніх літер посади, але
не особистий підпис, або розміщуватися на місці «М.~П.». В електронних листах
особистий підпис замінює кваліфікований електронний підпис або печатка
(\cl{4.4}).

\subsection{Відомості про виконавця (реквізит 27, \cl{5.27})}
Прізвище, власне ім'я виконавця та номер службового телефону:
«Федорчук Василь 123 45 67». Зазначаються в усіх вихідних листах, у лівому
нижньому куті лицьового боку останнього аркуша; за потреби — 2--3 виконавці.

\subsection{Інше}
Документи постійного та тривалого (понад 10 років) зберігання друкують на
одному боці аркуша. Реквізити бланка не підкреслюють і не відокремлюють
рискою (\cl{6.7}).

\subsection{Етикет ділового листування}
Звертання — у кличному відмінку: «Шановний пане Романе!», «Шановна пані
Олено!», «Шановні колеги!». Заключна формула «З повагою» доречна в ділових
листах, але не є реквізитом стандарту.

% =============================================================================
\section{Суміжні нормативні акти}

\begin{longtable}{@{}L{76mm}L{78mm}@{}}
\toprule
Акт & Значення для листів \\
\midrule
\endhead
ДСТУ 4163:2020 & склад, розміщення й оформлення реквізитів \\
Правила організації діловодства та архівного зберігання документів (наказ
  Мін'юсту від 18.06.2015 № 1000/5) & обов'язковість ДСТУ 4163 для юридичних осіб \\
Типова інструкція з діловодства (постанова КМУ від 17.01.2018 № 55) &
  діловодство в органах влади; з нею узгоджено форму «Власне ім'я ПРІЗВИЩЕ» \\
Правила надання послуг поштового зв'язку (постанова КМУ від 05.03.2009 № 270) &
  формат поштової адреси в реквізиті «Адресат» \\
Закон «Про звернення громадян», ст.~5 & звернення: прізвище, ім'я, по батькові,
  місце проживання, суть питання; письмове — з підписом і датою \\
Закон «Про доступ до публічної інформації», ст.~19--20 & запит: ім'я, поштова
  або електронна адреса, опис інформації, підпис і дата; відповідь — протягом
  5 робочих днів \\
Закон «Про забезпечення функціонування української мови як державної» &
  мова документів \\
\bottomrule
\end{longtable}

% =============================================================================
\section{Галерея прикладів}\label{sec:gallery}

\newcommand\galleryitem[3]{%
  \begin{minipage}[t]{0.31\linewidth}
    \centering\thumb[1]{#1}\par\vspace{1mm}
    {\small\textbf{#2}\par}
    {\footnotesize #3\par}
  \end{minipage}}

\noindent
\galleryitem{lyst-porozhnii.pdf}{templates/lyst-porozhnii}{Порожній шаблон з
  примітками: усі поля незаповнені, тому друкується лише текст.}\hfill
\galleryitem{lyst-povnyi.pdf}{templates/lyst-povnyi}{Усі можливості: емблема,
  гриф, два адресати, перелік додатків, два підписи, два виконавці.}\hfill
\galleryitem{00-blank-lysta.pdf}{00-blank-lysta}{Бланк для друку
  (\key{form = true}).}

\vspace{6mm}\noindent
\galleryitem{01-lyst-kutovyi-blank.pdf}{01-lyst-kutovyi-blank}{Лист-відповідь
  на кутовому бланку.}\hfill
\galleryitem{02-lyst-pozdovzhnii-blank.pdf}{02-lyst-pozdovzhnii-blank}{Поздовжній
  бланк підрозділу, два адресати, інтервал 1,5.}\hfill
\galleryitem{03-suprovidnyi-lyst.pdf}{03-suprovidnyi-lyst}{Супровідний лист
  із переліком додатків.}

\vspace{6mm}\noindent
\galleryitem{04-harantiinyi-lyst.pdf}{04-harantiinyi-lyst}{Гарантійний лист:
  назва виду документа, два підписи, «М.~П.».}\hfill
\galleryitem{05-dilovyi-lyst-fizosoby.pdf}{05-dilovyi-lyst-fizosoby}{Діловий
  лист фізичної особи (\key{sender = person}).}\hfill
\galleryitem{06-zvernennia-zaiava.pdf}{06-zvernennia-zaiava}{Заява громадянина
  (\key{blank = none}, \cs{applicant}).}

\vspace{6mm}\noindent
\galleryitem{07-zapyt-publichna-informatsiia.pdf}{07-zapyt-publichna-informatsiia}{Запит
  на публічну інформацію.}\hfill
\begin{minipage}[t]{0.31\linewidth}\end{minipage}\hfill
\begin{minipage}[t]{0.31\linewidth}\end{minipage}

% =============================================================================
\section{Поширені питання}

\textbf{«File `dstu4163.cls' not found».} Покладіть клас поруч із документом,
виконайте \texttt{make install} або компілюйте через \texttt{latexmk} у теках
пакета.

\textbf{Немає переносів, у журналі «Hyphen rules for 'ukrainian' set to
\textbackslash l@nil».} Встановіть шаблони: \verb|tlmgr install collection-langcyrillic|.

\textbf{У PDF не Times New Roman.} Шрифт не знайдено в системі; клас
використав Liberation Serif або TeX Gyre Termes. Встановіть Times New Roman
або задайте свій шрифт після \cs{documentclass}: \verb|\setmainfont{...}|.

\textbf{Довгий заголовок, посада чи адресат.} Ширина цих реквізитів
обмежена 73\,мм (\cl{7.6}); текст переноситься автоматично. Ручний розрив —
\verb|\\|.

\textbf{Більше чотирьох адресатів.} Стандарт вимагає список розсилки; клас
видає попередження, але друкує всіх.

\textbf{Чи можна використовувати Державний Герб?} Лише органам влади та
іншим особам, яким таке право надано законом. Іншим — емблема чи знак для
товарів і послуг (\key{emblem}).

\textbf{Як підписати лист електронно?} Скомпілюйте PDF і підпишіть його
кваліфікованим електронним підписом (КЕП) у сервісі на кшталт id.gov.ua чи
Дія. Розташування реквізитів лишається тим самим.

\textbf{Чим відрізняється проєкт стандарту?} Проєкт прДСТУ 4163:20\_\_ (2016)
містив інші значення; клас дотримується прийнятого ДСТУ 4163:2020.

\begin{longtable}{@{}L{56mm}L{44mm}L{52mm}@{}}
\toprule
Параметр & прДСТУ (2016) & ДСТУ 4163:2020 \\
\midrule
\endhead
Абзацний відступ & 12,5\,мм & 10\,мм \\
Адресат & 92\,мм & 90\,мм \\
Гриф затвердження / обмеження доступу & 104\,мм & 100\,мм \\
Розшифрування підпису & Ініціал(и), прізвище & Власне ім'я ПРІЗВИЩЕ \\
Виконавець & прізвище, ім'я, по батькові & Прізвище Власне ім'я \\
\bottomrule
\end{longtable}

% =============================================================================
\section{Історія версій}

\begin{description}[style=sameline,leftmargin=22mm,font=\normalfont\bfseries]
\item[1.1 (2026-11-18)] Умовні назви в усіх прикладах. Порожні реквізити й команди не друкуються, порожній
  бланк не залишає відступу; ключі \key{form}, \key{page-numbers},
  \key{label-*}; порожні рядки в \cs{dstusetup}; структура пакета
  (\texttt{templates}, \texttt{examples}, \texttt{doc}, \texttt{Makefile}),
  анотовані шаблони, ця документація.
\item[1.0 (2026-11-18)] Перша версія: кутовий, поздовжній бланки й аркуш
  без бланка; адресати, заголовок, додатки, підписи, виконавці; вісім
  прикладів.
\end{description}

% =============================================================================
\section{Ліцензія}

Цей твір можна поширювати та змінювати на умовах LaTeX Project Public
License версії 1.3c або (на ваш вибір) будь-якої пізнішої версії:
\url{https://www.latex-project.org/lppl.txt}.

Клас не є офіційним виданням стандарту. Для юридично значущих рішень
звіряйтеся з офіційним текстом ДСТУ~4163:2020 та інструкцією з діловодства
вашої установи.

\end{document}
